\documentclass[11pt,a4paper]{article}
\usepackage[utf8]{inputenc}
\usepackage[T1]{fontenc}
\usepackage[margin=2.5cm]{geometry}
\usepackage{lmodern}
\usepackage{parskip}
\usepackage{microtype}
\usepackage{booktabs}
\usepackage{longtable}
\usepackage{amsmath}
\usepackage{amssymb}
\usepackage{xcolor}
\usepackage{hyperref}
\usepackage{enumitem}
\usepackage{titlesec}
\usepackage{graphicx}

\hypersetup{
    colorlinks=true,
    linkcolor=blue!60!black,
    urlcolor=blue!60!black,
    citecolor=blue!60!black,
}

\titleformat{\section}{\large\bfseries}{}{0pt}{}[\titlerule]
\titleformat{\subsection}{\normalsize\bfseries}{}{0pt}{}
\titleformat{\subsubsection}{\normalsize\itshape}{}{0pt}{}

\title{\textbf{Harmonogramowanie na jednej maszynie ze wspólnym terminem realizacji}\\[4pt]
    \large Cel projektu, sformułowanie problemu i metody rozwiązania}
\author{Karolina Piotrowska, Jan Rosa}
\date{\today}

\begin{document}
\maketitle
\tableofcontents
\newpage

% ---------------------------------------------------------------------------
\section{Cel projektu}
% ---------------------------------------------------------------------------

Celem projektu \textbf{io-2026s} jest zbadanie, czy duże modele językowe (LLM)
mogą skutecznie automatyzować dobór hiperparametrów metaheurystyk stosowanych
do rozwiązywania kombinatorycznych problemów optymalizacyjnych.

Jako dziedzinę testową wybrano \emph{problem harmonogramowania na jednej
maszynie ze wspólnym terminem realizacji} (ang.\ \emph{common due date
single-machine scheduling}), ponieważ:
\begin{itemize}
  \item jest dobrze zbadany i posiada publicznie dostępne instancje
        benchmarkowe (biblioteka ORLib),
  \item jakość rozwiązania silnie zależy od właściwego doboru hiperparametrów
        metaheurystyki (co czyni go dobrym poligonem dla automatycznej
        optymalizacji hiperparametrów),
  \item instancje różnych rozmiarów ($n = 10, \ldots, 1000$) pozwalają ocenić
        skalowalność metod.
\end{itemize}

Projekt porównuje cztery podejścia:
\begin{enumerate}[label=\arabic*.]
  \item \textbf{Agenty bazowe} — losowy i zachłanny, stanowiące punkt odniesienia.
  \item \textbf{Klasyczne metaheurystyki} — wyżarzanie symulowane (SA)
        i algorytm genetyczny (GA) z ręcznie dobranymi hiperparametrami.
  \item \textbf{GEPA-SA / GEPA-GA} — te same metaheurystyki, lecz z
        hiperparametrami dobieranymi automatycznie przez framework GEPA
        sterowany przez LLM.
\end{enumerate}

% ---------------------------------------------------------------------------
\section{Sformułowanie problemu}
% ---------------------------------------------------------------------------

\subsection{Problem harmonogramowania ze wspólnym terminem realizacji}

Dane jest $n$ zadań, które muszą zostać wykonane na jednej maszynie
(bez możliwości przerywania).
Każde zadanie $i \in \{1, \ldots, n\}$ charakteryzuje się:

\begin{itemize}[noitemsep]
  \item $p_i > 0$ — deterministyczny czas przetwarzania,
  \item $a_i \geq 0$ — waga kary za \emph{zbyt wczesne} ukończenie,
  \item $b_i \geq 0$ — waga kary za \emph{zbyt późne} ukończenie.
\end{itemize}

Wszystkie zadania mają jeden wspólny termin realizacji:
\begin{equation}
  d = \left\lfloor h \cdot \sum_{i=1}^{n} p_i \right\rfloor,
  \label{eq:due_date}
\end{equation}
gdzie $h \in (0,1)$ jest \emph{współczynnikiem szczelności}.
Małe wartości $h$ (np.\ $0{,}2$) oznaczają bardzo napięty termin —
większość zadań jest spóźniona; duże wartości $h$ (np.\ $0{,}8$) dają
termin luźny — większość zadań kończy się zbyt wcześnie.

\subsection{Funkcja celu}

Szukamy permutacji $\pi$ zbioru zadań minimalizującej łączną ważoną karę
za wczesność i spóźnienie:
\begin{equation}
  \min_{\pi} \; C(\pi) = \sum_{i=1}^{n}
    \Bigl[ a_{\pi(i)} \cdot \max\!\bigl(0,\; d - C_{\pi(i)}\bigr)
         + b_{\pi(i)} \cdot \max\!\bigl(0,\; C_{\pi(i)} - d\bigr) \Bigr],
  \label{eq:objective}
\end{equation}
gdzie $C_{\pi(i)} = \sum_{j=1}^{i} p_{\pi(j)}$ jest czasem ukończenia
$i$-tego zadania w kolejności $\pi$.

Problem należy do klasy \textbf{NP-trudnych}~— dla dowolnej permutacji nie
istnieje wzór zamkniętej postaci optymalnego rozwiązania, wobec czego
konieczne jest przeszukiwanie przestrzeni permutacji.

\subsection{Instancje benchmarkowe (ORLib)}

Projekt korzysta z publicznie dostępnych instancji opracowanych przez
Biskupa i Feldmanna~\cite{biskup2001}:

\begin{center}
\begin{tabular}{@{}lrr@{}}
\toprule
\textbf{Plik} & \textbf{Liczba zadań} & \textbf{Rozmiar (KB)} \\
\midrule
\texttt{sch10}    &    10 & $<1$ \\
\texttt{sch20}    &    20 & $<1$ \\
\texttt{sch50}    &    50 &  $\sim$1 \\
\texttt{sch100}   &   100 &  $\sim$5 \\
\texttt{sch200}   &   200 & $\sim$15 \\
\texttt{sch500}   &   500 & $\sim$50 \\
\texttt{sch1000}  &  1000 & $\sim$200 \\
\bottomrule
\end{tabular}
\end{center}

Każdy plik zawiera 10 instancji problemowych.
Dla każdej instancji zdefiniowano 4 warianty terminu realizacji
($h \in \{0{,}2;\; 0{,}4;\; 0{,}6;\; 0{,}8\}$),
co daje łącznie 280 instancji benchmarkowych.
Najlepsze znane górne ograniczenia są podane w oryginalnej pracy.

% ---------------------------------------------------------------------------
\section{Opis rozwiązania}
% ---------------------------------------------------------------------------

\subsection{Środowisko i reprezentacja stanu}

Rozwiązanie jest reprezentowane jako permutacja
$\pi = (\pi(1), \pi(2), \ldots, \pi(n))$ zadań.
Agent operuje w środowisku \texttt{SchEnv} wzorowanym na Gym OpenAI:

\begin{description}[style=nextline, leftmargin=2.5cm]
  \item[\normalfont\textit{Stan}]
    Bieżąca permutacja zadań wraz ze skalarnymi sygnałami kontekstowymi:
    znormalizowany koszt bieżący, użyty ułamek budżetu kroków i współczynnik~$h$.
    Wektor obserwacji ma długość $3n + 3$.

  \item[\normalfont\textit{Akcja}]
    Wybór pary zadań $(i, j)$ z $i < j$ i zamiana ich pozycji w permutacji.
    Przestrzeń akcji ma rozmiar $\binom{n}{2}$.

  \item[\normalfont\textit{Nagroda}]
    \[
      r_t = \frac{C(\pi_t) - C(\pi_{t+1})}{C(\pi_0)},
    \]
    czyli znormalizowana poprawa kosztu w jednym kroku.
    Nagroda jest dodatnia, gdy zamiana poprawia harmonogram.
\end{description}

Agent wykonuje co najwyżej \texttt{max\_steps} kroków w jednym epizodzie.

\subsection{Inicjalizacja zachłanna}

Przed uruchomieniem przeszukiwania każdy agent inicjalizuje harmonogram
za pomocą prostej heurystyki zachłannej: zadania są sortowane rosnąco według
ilorazu $p_i / (a_i + b_i)$, co zapewnia sensowny punkt startowy.

% ---------------------------------------------------------------------------
\section{Metody rozwiązania}
% ---------------------------------------------------------------------------

\subsection{Agent losowy (Random)}

Jako dolny punkt odniesienia stosowany jest agent wybierający akcję
(parę do zamiany) jednostajnie losowo na każdym kroku.
Nie uczy się ani nie zapamiętuje historii; służy wyłącznie do
kwantyfikacji szumu w przestrzeni przeszukiwań.

\subsection{Agent zachłanny (Greedy)}

Na każdym kroku agent wyczerpująco sprawdza wszystkie $\binom{n}{2}$
możliwych zamian i wybiera tę, która przynosi największą natychmiastową
redukcję kosztu.
Złożoność jednego kroku wynosi $O(n^2)$, co czyni agenta powolnym dla
dużych $n$, ale gwarantuje monotoniczne zmniejszanie kosztu.

\subsection{Wyżarzanie symulowane (SA)}

Wyżarzanie symulowane przeszukuje przestrzeń
permutacji, akceptując ruchy pogarszające z prawdopodobieństwem zależnym
od temperatury.

\subsubsection*{Harmonogram temperatury}

Projekt stosuje \emph{oscylujący} harmonogram temperatury:
\begin{equation}
  T(t) = T_0 \cdot e^{-b \cdot t} + c \cdot \sin^2\!\left(\frac{t}{d}\right),
  \label{eq:temperature}
\end{equation}
gdzie:
\begin{itemize}[noitemsep]
  \item $T_0$ — temperatura początkowa,
  \item $b$ — współczynnik wykładniczego zaniku,
  \item $c$ — amplituda periodycznego podgrzewania (sinusoidalnego),
  \item $d$ — okres podgrzewania (w krokach),
  \item $T_{\min}$ — dolne ograniczenie temperatury.
\end{itemize}

Ruch do sąsiedniego harmonogramu $\pi'$ (różniącego się jedną zamianą) jest
akceptowany zawsze, gdy $C(\pi') \leq C(\pi)$, albo z prawdopodobieństwem
$\exp\!\bigl(-\Delta C / T(t)\bigr)$ w przeciwnym razie.

Sinusoidalny człon $c \cdot \sin^2(t/d)$ wprowadza periodyczne skoki
temperatury (podgrzewanie), które pomagają agentowi uciekać z lokalnych
minimów.

\subsubsection*{Hiperparametry SA}

\begin{center}
\begin{tabular}{@{}lll@{}}
\toprule
\textbf{Parametr} & \textbf{Zakres} & \textbf{Wpływ} \\
\midrule
$T_0$      & $[1, 10000]$   & Swoboda eksploracji na początku \\
$b$        & $[10^{-6}, 0{,}1]$ & Szybkość chłodzenia \\
$c$        & $[0, 100]$     & Intensywność podgrzewania \\
$d$        & $[1, 500]$     & Częstotliwość podgrzewania \\
$T_{\min}$ & $[10^{-8}, 1]$ & Minimalna temperatura \\
\bottomrule
\end{tabular}
\end{center}

\subsection{Algorytm genetyczny (GA)}

Algorytm genetyczny utrzymuje populację kandydatów
(permutacji) i ewoluuje ją przez selekcję, krzyżowanie i mutację.

\subsubsection*{Operatory genetyczne}

\begin{description}[style=nextline]
  \item[\normalfont\textit{Krzyżowanie (Order Crossover, OX)}]
    Wybierane jest losowe okno w permutacji rodzica~1; geny z okna są
    kopiowane bezpośrednio do potomka.
    Pozostałe pozycje uzupełniane są zadaniami z rodzica~2 w kolejności ich
    występowania, pomijając już wstawione.
    OX zachowuje względną kolejność zadań i nie tworzy zduplikowanych genów.
    Operacja krzyżowania całego pokolenia jest wykonywana w jednym wywołaniu
    rozszerzenia C (\texttt{crossover\_opt.c}), co znacznie skraca czas
    obliczeń.

  \item[\normalfont\textit{Mutacja}]
    Z prawdopodobieństwem \texttt{mutation\_rate} losowa para pozycji
    w potomku jest zamieniana miejscami.

  \item[\normalfont\textit{Elityzm}]
    Najlepszych \texttt{elite\_size} osobników jest kopiowanych do kolejnego
    pokolenia bez zmian.
\end{description}

\subsubsection*{Osadzenie GA w epizodzie}

GA nie działa niezależnie, lecz jako \emph{wewnętrzna pętla optymalizacji}
uruchamiana co \texttt{reoptimize\_every} kroków epizodu.
Całkowita liczba pokoleń GA wynosi zatem:
\[
  \text{pokolenia łącznie} = \left\lfloor \frac{\texttt{max\_steps}}
  {\texttt{reoptimize\_every}} \right\rfloor \times \texttt{generations}.
\]

\subsubsection*{Hiperparametry GA}

\begin{center}
\begin{tabular}{@{}lll@{}}
\toprule
\textbf{Parametr} & \textbf{Zakres} & \textbf{Wpływ} \\
\midrule
\texttt{population\_size}   & $[10, 500]$  & Szerokość przeszukiwania \\
\texttt{generations}        & $[1, 100]$   & Głębokość przeszukiwania \\
\texttt{mutation\_rate}     & $[0, 1]$     & Różnorodność populacji \\
\texttt{elite\_size}        & $[1, 20]$    & Zachowanie najlepszych \\
\texttt{reoptimize\_every}  & $[1, 20]$    & Częstość uruchamiania GA \\
\bottomrule
\end{tabular}
\end{center}

\subsection{GEPA — automatyczna optymalizacja hiperparametrów przez LLM}

\subsubsection*{Motywacja}

Wyniki SA i GA są silnie uzależnione od wyboru hiperparametrów.
Ręczne strojenie jest kosztowne i wymaga ekspertyzy dziedzinowej.
Framework GEPA zastępuje ten proces pętlą refleksji LLM,
w której model językowy analizuje wyniki ewaluacji i proponuje ulepszoną
konfigurację.

\subsubsection*{Pętla optymalizacji GEPA}

\begin{enumerate}
  \item \textbf{Inicjalizacja:} ziarno konfiguracji (domyślne hiperparametry)
        jest przekazywane do GEPA.
  \item \textbf{Ewaluacja bieżącej konfiguracji:} agent jest uruchamiany
        na losowym podzbiorze instancji treningowych; wynik jest zapisywany.
  \item \textbf{Refleksja LLM:} bieżąca konfiguracja, wyniki ewaluacji
        (poprawa\%, wzorzec zbieżności, czas) oraz prompt dziedzinowy
        są przekazywane do LLM.  Model odpowiada obiektem JSON z propozycją
        nowych wartości hiperparametrów.
  \item \textbf{Akceptacja / odrzucenie:} nowa konfiguracja jest ewaluowana
        na tym samym podzbiorze i porównywana z bieżącą.
        Jeśli wynik złożony (jakość $-$ kara za czas) jest wyższy,
        konfiguracja zostaje zaakceptowana.
  \item Kroki 2--4 są powtarzane aż do wyczerpania budżetu
        (\texttt{max\_metric\_calls}).
  \item \textbf{Walidacja:} najlepsza znaleziona konfiguracja jest
        ewaluowana na pełnym zbiorze walidacyjnym.
\end{enumerate}

\subsubsection*{Funkcja oceny (score)}

Ocena złożona uwzględnia zarówno jakość rozwiązania, jak i czas wykonania:
\begin{equation}
  \text{score} = \underbrace{\frac{C_0 - C^*}{C_0}}_{\text{jakość}}
               - \alpha \cdot \underbrace{\frac{t}{t_{\max}}}_{\text{norma czasu}},
  \label{eq:score}
\end{equation}
gdzie $C_0$ to koszt rozwiązania zachłannego, $C^*$ to najlepszy znaleziony
koszt, $t$ to czas jednego uruchomienia agenta, $t_{\max}$ to maksymalny
czas obserwowany dotychczas, a $\alpha = 0{,}2$ to waga kary za czas.

\subsubsection*{Modele LLM}

GEPA komunikuje się z modelem przez interfejs LiteLLM.
Projekt obsługuje modele Ollama uruchamiane lokalnie:

\begin{center}
\begin{tabular}{@{}ll@{}}
\toprule
\textbf{Model} & \textbf{Klucz CLI} \\
\midrule
Qwen3 4B Instruct  & \texttt{ollama/qwen3:4b-instruct-2507-q4\_K\_M} \\
Gemma4 2B          & \texttt{ollama/gemma4:e2b} \\
Phi4-mini          & \texttt{ollama/phi4-mini-reasoning:latest} \\
Granite4 1B        & \texttt{ollama/granite4:1b-h-q8\_0} \\
\bottomrule
\end{tabular}
\end{center}

\subsubsection*{Prompt refleksji}

Prompt dostarczany modelowi zawiera:
\begin{enumerate}[noitemsep]
  \item Opis semantyczny hiperparametrów i ich zakresy dopuszczalnych wartości.
  \item Wskazówki dotyczące interpretacji wyników (poprawa\%, wzorzec zbieżności,
        stosunek liczby korzystnych ruchów do wszystkich kroków).
  \item Strategię eksploracji — model jest instruowany, by proponować
        \emph{istotnie różne} konfiguracje od bieżącej i nie powtarzać
        wcześniej odrzuconych.
  \item Bieżącą konfigurację hiperparametrów i wyniki ewaluacji na instancjach
        treningowych.
\end{enumerate}

Model odpowiada \emph{wyłącznie} obiektem JSON z nowymi wartościami
hiperparametrów.  Odpowiedzi nieparsowalne są zastępowane losową konfiguracją
z dopuszczalnego zakresu, co pozwala kontynuować optymalizację nawet
w przypadku słabszych modeli.

% ---------------------------------------------------------------------------
\section{Optymalizacje implementacyjne}
% ---------------------------------------------------------------------------

\subsection{Rozszerzenia C}

Krytyczna ścieżka obliczeniowa (ewaluacja kosztu i krzyżowanie OX) jest
zrealizowana jako rozszerzenia Pythona skompilowane z C:

\begin{description}[style=nextline]
  \item[\texttt{evaluate\_opt.c}]
    Ewaluacja kosztu permutacji i przyrostowa kalkulacja $\Delta C$ dla
    pojedynczej zamiany w jednym wywołaniu C, bez alokacji Pythona.
  \item[\texttt{crossover\_opt.c}]
    Krzyżowanie OX całego pokolenia ($k \times n$) w jednym wywołaniu C,
    zwracające tablicę \texttt{int64} bez pośrednictwa Pythona.
\end{description}

\subsection{Redukcja narzutu oceny}

Podczas treningu GEPA każda konfiguracja jest ewaluowana \textbf{3-krotnie}
i uśredniana, co redukuje wpływ szumu stochastycznego bez nadmiernego
wydłużania czasu obliczeń.

\begin{thebibliography}{1}

\bibitem{biskup2001}
D.~Biskup, M.~Feldmann,
\textit{Benchmarks for scheduling on a single machine against restrictive and
unrestrictive common due dates},
Computers \& Operations Research, vol.~28, pp.~787--801, 2001.

\end{thebibliography}

\end{document}
